% !TEX root = task-3-20-3.tex
% Editable task bank. Append new BankTask blocks at the end of the document.
\documentclass[12pt,a4paper]{article}
\usepackage{physics-textbook}
\usepackage{graphicx,needspace}
\usetikzlibrary{decorations.pathmorphing}
\newif\ifShowSolutions
\ifdefined\TutorEdition
  \ShowSolutionstrue
\else
  \ShowSolutionsfalse % Change to \ShowSolutionstrue to display solutions.
\fi
\newcounter{banktask}
\newsavebox{\BankTaskBox}
\newcommand{\BankTaskContent}[3]{%
  \noindent\textbf{\thebanktask.}\quad #2\par
  \ifShowSolutions
    \Needspace{50mm}% Keep the solution heading with its opening text or diagram.
    #3\par
    \ifdefined\TutorEdition
      \ifcsname TutorSource#1\endcsname
        \nopagebreak\smallskip{\small\csname TutorSource#1\endcsname\par}%
      \fi
    \fi
  \fi
}
\newcommand{\BankTask}[3]{%
  \refstepcounter{banktask}\par\medskip
  \sbox{\BankTaskBox}{\begin{minipage}{\linewidth}%
    \BankTaskContent{#1}{#2}{#3}\end{minipage}}%
  \ifdim\dimexpr\ht\BankTaskBox+\dp\BankTaskBox\relax>\dimexpr\textheight-20mm\relax
    \Needspace{35mm}\BankTaskContent{#1}{#2}{#3}%
  \else
    \Needspace{\dimexpr\ht\BankTaskBox+\dp\BankTaskBox\relax}%
    \noindent\usebox{\BankTaskBox}\par
  \fi
  \medskip
}
\begin{document}
\BookSubsection{3.20.3}{Энергия магнитного поля}


% Source 458c72: https://ege.fipi.ru/bank/questions.php?search=1&proj=BA1F39653304A5B041B656915DC36B38&qid=458c72
\BankTask{458c72}{
Чему равна индуктивность катушки (в $\text{мГн}$), если при силе тока \(5\,\text{А}\) энергия её магнитного поля равна \(12{,}5\,\text{мДж}\)?
}{
\smallskip\noindent\textbf{Решение от GPT}\par
Энергия магнитного поля катушки \(W=\dfrac{LI^2}{2}\), поэтому
\[
L=\frac{2W}{I^2}=\frac{2\cdot12{,}5\cdot10^{-3}}{5^2}=10^{-3}\,\text{Гн}=1\,\text{мГн}.
\]
Ответ: \(1\).
}

% Source D747BC: https://ege.fipi.ru/bank/questions.php?search=1&proj=BA1F39653304A5B041B656915DC36B38&qid=D747BC
\BankTask{D747BC}{
Какова энергия магнитного поля катушки индуктивностью \(0{,}4\,\text{Гн}\), если сила тока в катушке равна \(6\,\text{А}\) (в $\text{Дж}$)?
}{
\smallskip\noindent\textbf{Решение от GPT}\par
Энергия магнитного поля катушки:
\[
W=\frac{LI^2}{2}=\frac{0{,}4\cdot 6^2}{2}=7{,}2\,\text{Дж}.
\]
Ответ: \(7{,}2\,\text{Дж}\).
}

% Source C453EB: https://ege.fipi.ru/bank/questions.php?search=1&proj=BA1F39653304A5B041B656915DC36B38&qid=C453EB
\BankTask{C453EB}{
Энергия магнитного поля катушки с током равна \(0{,}64\,\text{Дж}\). Индуктивность катушки равна \(20\,\text{мГн}\). Какова сила тока в катушке (в $\text{А}$)?
}{
\smallskip\noindent\textbf{Решение от GPT}\par
Из \(W=LI^2/2\), учитывая \(L=20\,\text{мГн}=0{,}020\,\text{Гн}\), получаем:
\[
I=\sqrt{\frac{2W}{L}}=\sqrt{\frac{2\cdot0{,}64}{0{,}020}}=8\,\text{А}.
\]
Ответ: \(8\,\text{А}\).
}

% Source 9CFC20: https://ege.fipi.ru/bank/questions.php?search=1&proj=BA1F39653304A5B041B656915DC36B38&qid=9CFC20
\BankTask{9CFC20}{
В катушке индуктивностью $20\,\text{мГн}$ сила тока $I$ зависит от времени $t$, как показано на графике, приведённом на рисунке.

\begin{enumerate}
\item Модуль ЭДС самоиндукции в катушке максимален в интервале времени от $0$ до $1\,\text{с}$.
\item Энергия магнитного поля катушки в интервале времени от $2$ до $3\,\text{с}$ равна $40\,\text{мДж}$.
\item Модуль ЭДС самоиндукции в катушке в интервале времени от $3$ до $5\,\text{с}$ равен $10\,\text{мВ}$.
\item Модуль ЭДС самоиндукции в катушке минимален в интервале времени от $3$ до $4\,\text{с}$.
\item Модуль скорости изменения тока в катушке максимален в интервале времени от $5$ до $6\,\text{с}$.
\end{enumerate}

\begin{center}
\begin{tikzpicture}[book diagram,x=0.75cm,y=0.75cm]
\draw[thin] (0,1)--(6,1) (0,2)--(6,2) (0,3)--(6,3) (0,4)--(6,4);
\foreach \x in {1,2,3,4,5,6} {\draw[thin] (\x,0)--(\x,4);}
\draw[-{book tip}] (0,0)--(6.75,0) node[right] {$t,\,\text{с}$};
\draw[-{book tip}] (0,0)--(0,4.9) node[above left] {$I,\,\text{А}$};
\node[below left] at (0,0) {$0$};
\foreach \x in {1,2,3,4,5,6} {\node[below] at (\x,0) {$\x$};}
\foreach \y in {1,2,3,4} {\node[left] at (0,\y) {$\y$};}
\draw[thick] (0,0)--(2,4)--(3,4)--(5,3)--(6,0);
\end{tikzpicture}
\end{center}
}{
\smallskip\noindent\textbf{Решение от GPT}\par
Используем $|\mathcal E|=L|\Delta I/\Delta t|$: на интервалах $0$–$2$, $2$–$3$, $3$–$5$ и $5$–$6\,\text{с}$ модуль скорости изменения тока равен соответственно $2$, $0$, $0{,}5$ и $3\,\text{А}/\text{с}$, поэтому утверждение 1 неверно, а 5 верно. При $2<t<3\,\text{с}$ энергия $W=LI^2/2=0{,}02\cdot4^2/2=0{,}16\,\text{Дж}=160\,\text{мДж}$, поэтому утверждение 2 неверно. На интервале $3$–$5\,\text{с}$ получаем $|\mathcal E|=0{,}02\cdot0{,}5=0{,}01\,\text{В}=10\,\text{мВ}$, поэтому утверждение 3 верно. Минимальная ЭДС равна нулю на интервале $2$–$3\,\text{с}$, поэтому утверждение 4 неверно. Ответ: 3, 5.
}

\end{document}
